\documentclass[11pt]{article}
\usepackage{amsmath,amssymb,mathtools}
\usepackage{fontspec}
\usepackage{unicode-math}
\setmainfont{Latin Modern Roman}
\setmathfont{Latin Modern Math}
\usepackage{xurl}
\usepackage[unicode,colorlinks=true,linkcolor=blue,urlcolor=blue,pdftitle={Problem statement}]{hyperref}
\newcommand{\sref}[2]{\hyperlink{source:#1:#2}{[S#2]}}
\newcommand{\cataloguescope}[1]{\paragraph{Mathematical statement.}#1}
\begin{document}
\section{Research Problems in Function Theory — Problem 2.24}\label{statement}

\subsection{Definitions and mathematical statement}

\paragraph{Shared notation.}$\mathbb N=\{1,2,\ldots\}$, and $\mathbb Z,\mathbb Q,\mathbb R,\mathbb C$ denote the integers, rationals, reals and complex numbers. Any other conventions are specified locally.


For an entire function $f$, put $Z(f)=\{z\in\mathbb C:f(z)=0\}$ and $O(f)=\{z\in\mathbb C:f(z)=1\}$.

\paragraph{Statement recorded in the supplied source.}
Can there exist an entire function $f$ and two distinct straight lines $\ell,m$ such that
\[Z(f)\cup O(f)\subseteq\ell\cup m,\]
and the combined set $Z(f)\cup O(f)$ has infinitely many points on each of $\ell$ and $m$? \sref{statement}{2}
The two-line condition refers to their union: it does not require all zeros to lie on one line and all ones on the other. Update 2.24 makes this distinction explicit through the example $f(z)=\sin(z^2)$, whose zero-points and one-points all lie on the real or imaginary axis. The original paragraph separately discusses Edrei's obstruction to assigning the two levels to two intersecting lines. \sref{statement}{2}
\subsection{Short English statement}
Can all zero-points and one-points of an entire function be confined to two lines, with infinitely many level-set points on each line?
\subsection{Sources}
\begin{description}
\item[\hypertarget{source:statement:1}{S1}] ULAM.AI and the UnsolvedMath Contributors, \emph{UnsolvedMath: A Curated Collection of Open Mathematics Problems}, catalogue entry. CC BY 4.0. \url{https://huggingface.co/datasets/ulamai/UnsolvedMath}.
\item[\hypertarget{source:statement:2}{S2}] W. K. Hayman and E. F. Lingham, \emph{Research Problems in Function Theory} (2018), Chapter 2, Problem 2.24 and Update 2.24. \url{https://arxiv.org/abs/1809.07200}.
\end{description}
\label{end:statement}
\end{document}
